\chapter{Information Theory and Linguistic Stability}

\section{Stochastic Modeling of Linguistic Evolution}
Linguistic change is characterized by entropic decay, where high-resolution structural and lexical information is eroded over temporal horizons. This study utilizes an information-theoretic framework to model the tiered stability of phonological and semantic signals, identifying specific 'conserved regions' that persist over deep-time horizons.

\section{The Decay Threshold and Signal Horizons}
Classic glottochronological models utilize a uniform decay rate (approximately 14\% per millennium), which suggests a limit of recognizable signal recovery at roughly 10,000 BP. We hypothesize that the recovery of signals at the 40,000 BP (PED) horizon is contingent upon isolating the ultra-conservative structural invariants.

\section{Differential Stability Analysis and Frequency of Use}
We identify a non-uniform distribution of lexical stability across word classes, heavily governed by the frequency-of-use hypothesis (e.g., Pagel et al., 2013). The frequency with which a linguistic item is used directly correlates with its resistance to lexical replacement.
\begin{itemize}
    \item \textbf{Dynamic Strata:} Nouns representing technological artifacts or culturally specific flora/fauna exhibit high replacement rates (15–20\% per millennium).
    \item \textbf{Conserved Strata:} Core elements, specifically deictic pronouns ('I', 'thou') and interrogatives, which are utilized at frequencies exceeding 1 in 1,000 spoken words.
\end{itemize}

\subsection{Monte Carlo Validation: The 34,000-Year Half-Life}
Monte Carlo simulations were executed to model the survival of the Deictic core over 40,000 years. The underlying exponential decay equation, $S = S_0(1 - r)^t$, yields an expected retention of $0.4457$ (approximately 44 words out of 100) when the decay rate $r = 0.02$ (2\% per millennium) and $t = 40$. 

\begin{table}[h]
\centering
\begin{tabular}{lcc} \toprule
\textbf{Word Class} & \textbf{Decay Rate (/1ky)} & \textbf{Survival Prob. (40ky)} \\ \midrule
Standard Nouns & 15.0\% & 0.15\% \\
Deictic Core & 2.0\% & 44.6\% \\ \bottomrule
\end{tabular}
\caption{Survival Thresholds in Deep-Time Reconstruction.}
\end{table}

A 2\% decay rate implies a lexical half-life of roughly 34,000 years. While this represents extreme stability, it is entirely consistent with empirical studies of ultra-conserved words in Indo-European and Dravidian. The neurological entrenchment of basic deictic reference frames means their replacement requires a catastrophic collapse of the communicative network. Thus, the PED "Deictic Core" defied standard entropy not through geographic isolation, but through sheer frequency of use and indispensable structural function. The signal recovered here is the mathematical residue of the most frequently spoken syllables of the Upper Paleolithic.
